\documentclass[10pt,a4paper]{amsart}
\usepackage[margin=19mm]{geometry}
\usepackage{amsmath,amssymb,mathtools}
\usepackage[scr=rsfs,cal=euler]{mathalfa}
\usepackage{xfrac}
\usepackage[unicode,hidelinks,bookmarks=false]{hyperref}
\newcommand{\ie}{i.e.\ }
\newcommand{\cf}{cf.\ }
\newcommand{\resp}{respectively\ }
\newcommand{\Subsection}{\subsection}
\providecommand{\nobreakdash}{\nobreak\hbox{-}}
\setlength{\emergencystretch}{2em}
\multlinegap0pt
\usepackage{amssymb}
\usepackage{dsfont}
\usepackage{float}
\newtheorem{lemma}{Lemma}
\title{On the Kobayashi isometries of a class of 2-dimensional Lempert manifolds}
\author{Naveen Gupta}
\date{Source: arXiv:2609.04251v1 (CC BY 4.0)}
\begin{document}
\maketitle

\noindent\emph{Source and licence.} On the Kobayashi isometries of a class of 2-dimensional Lempert manifolds. Version arXiv:2609.04251v1 (CC BY 4.0), distributed by arXiv under the Creative Commons Attribution 4.0 International licence, http://creativecommons.org/licenses/by/4.0/, as recorded in the arXiv licence metadata for this exact version and shown on the abstract page https://arxiv.org/abs/2609.04251v1. Full licence notice: \texttt{licenses/arxiv-2609.04251-CC-BY-4.0.txt}.\par\smallskip

\noindent\textit{Editorial scope.} Counted unit: the article's lemma lem:connect (source0.tex lines 132--136). The article's complete original proof of it is delivered with it (source0.tex lines 138--152, one \texttt{proof} environment of 15 lines). No mathematics is written, repaired or re-derived: the statement and the proof are the article's own lines, in the article's own notation, and no retained line is substituted or re-flowed.  No further source-local context is needed: every cross-reference of every retained line resolves inside the two delivered spans.  Nothing was omitted from any retained span, so the package carries no omission record.  No retained line contains a citation, so the wrapper carries no bibliography and no bibliography entry.  No retained line contains a figure environment, an image-inclusion command or drawing code, so no image is delivered and the package declares no asset dependency.  Editorial formatting only, in addition: an AMS \texttt{amsart} preamble that restates the article's own declarations and macros used by the retained lines, namely the environments \texttt{lemma} on separate counters, together with the standard packages those lines need.  The retained environments print in this document as Lemma 1 in order of appearance, whereas the article prints the same items with the numbering of its own full document.  The complete native source of the article as distributed in the arXiv e-print for this exact version is bundled unchanged as \texttt{sources/209353/source0.tex}.
\begin{lemma}\label{lem:connect}
Let $M$ be a Lempert manifold whose Carath\`{e}odory universal set 
$\mathcal{F}=\{f_1,f_2,\ldots,f_m\}$ is finite. Then the  Carath\`{e}odory-Rieffen 
metric $\gamma_M:TM\to [0,\infty)$ is $C_{\text{loc}}^{0,1}$.
\end{lemma}

\begin{proof}
	Recall that  for any $p\in M$ and $X\in T_pM$, we have 
	$$\gamma_M(p;X)=\sup_{f\in \mathcal{F}}\gamma_{\mathbb{D}}(f(p);f'(p)X).$$
	
	Since the Carath\`{e}odory universal set is finite, the expression reduces to 
		$$\gamma_M(p;X)=\max_{1\leq i\leq m}\gamma_{\mathbb{D}}(f_i(p);f_i'(p)X).$$
		
		For any fixed $1\leq i\leq m$,  consider the map 
		\begin{align*}(p,X)\to& \gamma_{\mathbb{D}}(f_i(p);f_i'(p)X)\\
			& =\dfrac{|f_i'(p)X|}{1-|f_i(p)|}.			
		\end{align*}
	Since $f_i$ is holomorphic and the image of $f_i$ lies within $\mathbb{D}$, it follows that 
	this map is locally Lipschitz. Therefore, $\gamma_M$, which is maximum of these Lipschitz 
	functions, is also Lipschitz. This proves that $\gamma_M$ is $C_{\text{loc}}^{0,1}$.
\end{proof}

\end{document}
